\documentclass[11pt,twocolumn]{article}
\usepackage[utf8]{inputenc}
\usepackage{amsmath}
\usepackage{amssymb}
\usepackage{graphicx}
\usepackage{booktabs}
\usepackage{microtype}

\title{\textbf{The Rhythmic Tesseract: Multi-Dimensional Spatial Mapping and Permutation Engines in Complex Polymetric Audio Architectures}}
\author{\small\textit{Audio Engine Architecture and Rhythmic Theory Research Group} \\
\small Internal Technical Specification Monograph}
\date{May 2026}

\begin{document}

\maketitle

\begin{abstract}
This paper presents a formal mathematical model for tracking multi-layered polymetric structures by mapping linear time signatures to three-dimensional geometric vectors. By nesting an independent polyrhythmic matrix (the inner cube) within a rigid master frame (the outer cube), we eliminate cumulative clock drift in high-priority audio threads. Furthermore, we demonstrate that structural permutations function as rotational operators within a discrete phase space, creating a dynamic, self-correcting temporal lattice.
\end{abstract}

\section{Introduction}
Typical digital audio workstation (DAW) metronomes treat temporal synchronization as a flat, linear timeline. While sufficient for symmetric time signatures, this approach breaks down under complex polymetric conditions where independent cyclic structures (e.g., $6/8, 3/4, 4/4$) phase simultaneously over an asymmetric baseline ($5/4$). 

To manage this without compounding floating-point jitter, we formalize a spatial model—inspired by the geometric and polyphonic drum patterns popularized by Carey \cite{drumeo_carey}, the structural binary change-matrices analyzed historically by Leibniz \cite{leibniz_binary}, and the systematic permutation logic found in traditional scale classification systems \cite{venkatasubramanian_melakarta}—wherein time signatures fold into an integrated three-dimensional manifold.

\section{The Coordinate System}
We define the master audio clock as a three-dimensional vector space $\mathbb{R}^3$, where any discrete audio event or transient impact point $P$ is determined by the coordinate vector:
\begin{equation}
P = \begin{bmatrix} x \\ y \\ z \end{bmatrix}
\end{equation}

The axes map explicitly to the structural parameters of the audio engine:
\begin{itemize}
    \item \textbf{X-Axis (Linear Width):} The macro-tempo or global horizontal timeline, governed by the Beats Per Minute (BPM) parameter.
    \item \textbf{Y-Axis (Vertical Density):} The concurrent layer depth, tracking stacked metric layers.
    \item \textbf{Z-Axis (Subdivision Depth):} The discrete micro-slice phase, determined by the hyper-fine tick resolution.
\end{itemize}

\section{The Nested Cube Manifold}
The system isolates global boundaries from local shifting ostinatos by maintaining two separate, concentric geometric frames sharing a single absolute origin point $\mathcal{O}(0,0,0)$.

\subsection{The Outer Cube (Master Frame)}
The outer container represents the foundational metric boundary (e.g., $5/4$ time). It establishes the master buffer constraints and dictates the absolute downbeat boundary $B_{\text{master}}$. The volume $V_{\text{outer}}$ of this structural frame is bounded by:
\begin{equation}
V_{\text{outer}} = \{ (x,y,z) \in \mathbb{R}^3 \mid 0 \le x \le M_{\text{beats}} \}
\end{equation}
where $M_{\text{beats}} = 5$ in a standard asymmetric pentagonal baseline.

\subsection{The Inner Cube (Polyrhythmic Matrix)}
Suspended within $V_{\text{outer}}$ is the inner matrix generating the sub-meters ($6/8, 3/4, 4/4$). Because these inner metrics do not share divisor symmetry with the outer frame, they are mathematically modeled as a secondary coordinate system rotating relative to the primary container.

The phase alignment of any inner transient point $P_{\text{inner}}$ relative to the outer container is determined by computing its spatial vector distance from the shared origin:
\begin{equation}
D(P) = \sqrt{(x_{\text{out}} - x_{\text{in}})^2 + (y_{\text{out}} - y_{\text{in}})^2 + (z_{\text{out}} - z_{\text{in}})^2}
\end{equation}

By resolving all threads to the absolute origin $\mathcal{O}(0,0,0)$ on every hardware callback tick, cumulative pointer drift is forced to absolute zero.

\section{Permutations as Rotational Operators}
Permutations within the inner ostinatos do not merely alter accent placements; they function as discrete rotational matrices that tilt the inner cube along the Z-axis (Depth). 

Let a rhythmic cell of length $N$ micro-slices be represented as a discrete state sequence. A metric permutation shift of $\sigma$ slices corresponds to a rotation operator $\mathbf{R}_z(\theta)$ acting on the phase space, defined by:
\begin{equation}
\mathbf{R}_z(\theta) = \begin{bmatrix} 
\cos\theta & -\sin\theta & 0 \\ 
\sin\theta & \cos\theta & 0 \\ 
0 & 0 & 1 
\end{bmatrix}
\end{equation}
where the phase angle displacement $\theta$ is a function of the discrete permutation offset:
\begin{equation}
\theta = 2\pi \left( \frac{\sigma}{N_{\text{slices}}} \right)
\end{equation}

\begin{table}[h]
\centering
\small
\caption{\small Geometric Manifestation of Rhythmic Shifts}
\label{tab:permutations}
\vspace{2mm}
\begin{tabular}{p{2.2cm} p{2.8cm} p{2.4cm}}
\toprule
\textbf{Shift ($\sigma$)} & \textbf{Geometric Translation} & \textbf{Engine State} \\ 
\midrule
$\sigma = 0$               & Uniform Inner Alignment        & Root Phase \\ \addlinespace[0.5em]
$\sigma = +1$              & Forward Angular Tilt           & Displaced Accent \\ \addlinespace[0.5em]
$\sigma = +2$              & Diagonal Core Torsion          & Max Phase Drift \\ 
\bottomrule
\end{tabular}
\end{table}

When $\sigma \ne 0$, the internal corners of the inner cube spin entirely out of parallel alignment with the outer walls (Table \ref{tab:permutations}). The structural integrity of the groove is maintained exclusively by the hyper-fine micro-slices acting as tension struts between the two frames, until the macro-cycle completes and both cubes snap back to $\mathcal{O}(0,0,0)$.

\begin{thebibliography}{9}
\bibitem{drumeo_carey}
Drumeo, \emph{Tool Drummer Danny Carey: Reasons Why He's A Rhythmic Genius}, Drumeo Editorial Network, 2023.

\bibitem{leibniz_binary}
G. W. Leibniz, \emph{Explication de l'Arithmétique Binaire (Explanation of Binary Arithmetic)}, Histoires de l'Académie Royale des Sciences, Paris, 1703.

\bibitem{venkatasubramanian_melakarta}
R. Venkatasubramanian, \emph{Mathematics of Melakartha Ragas in Carnatic Music}, Center for Advanced Computation and Telecommunications, 2014.
\end{thebibliography}

\end{document}
